If we wanted to design tests that could differentiate genuine curiosity from simulation, the literature suggests several candidates:

\begin{enumerate}[leftmargin=*]
  \item \textbf{Stress tolerance under genuine uncertainty.} Does the system explore \emph{more} when stakes are high and outcomes uncertain? Humans with high curiosity do. Current LLMs do the opposite \citep{wang2025}.

  \item \textbf{Curiosity about one's own processes.} Metacognitive curiosity---being curious about \emph{why you are curious}---is a distinctly human pattern. If an AI system spontaneously investigates its own exploratory drives without being prompted, that is at least interesting.

  \item \textbf{Persistence without reward signal.} Human curiosity persists even when information-seeking is costly and unrewarded. Harlow's monkeys solved puzzles with no food reward. Does an AI system continue exploring when the exploration bonus is removed?

  \item \textbf{Novel question generation.} Not recombination of known question templates, but genuinely new questions that could not have been predicted from training data. The distinction between ``analytical follow-ups'' and ``a question that slaps you in the forehead'' maps onto this---whims without basis, questions that surprise even the asker.

  \item \textbf{Developmental trajectory.} If an AI system shows genuine curiosity, it should show something like developmental change---early broad exploration giving way to deepening expertise, with periodic bursts of novel exploration. The U-shaped curve.
\end{enumerate}
